\begin{corollary}[Intrinsic bit law]\label{cor:intrinsic-bits}
Let $B_*(\varepsilon,a)$ be the least integer persistent bit budget
for maximum-history streaming regret at most $\varepsilon$ in this
experiment. Uniformly for $0<\varepsilon\le1$ and $a\in K$,
\begin{equation}\label{eq:intrinsic-bits}
 B_*(\varepsilon,a)=
 \max_{\substack{1\le n<N\\1\le\ell\le p_{n,N-n}}}
 \left[\frac{\ell}{2}\log_2\frac1\varepsilon+
       \log_2\mathcal V_{N-n,\ell}(a)\right]_++O(1).
\end{equation}
A zero determinant volume contributes zero to this maximum, by interpreting
its logarithm as $-\infty$ before taking the positive part. The same law
holds for the unconditional maximum-checkpoint criterion. At each fixed
calibration the small-error slope is $D_A(N)/2$.
\end{corollary}
\begin{proof}
With $M=2^B$, solving each inequality
$(\mathcal V_{m,\ell}/M)^{2/\ell}\le\varepsilon$ gives the displayed
lower bound on $B$. Constants in Theorem~\ref{thm:resolution-main} change
it by at most a bounded amount, since $\ell$ is bounded. Integer rounding
costs at most one further bit. For fixed $a$, the largest nonzero index
at checkpoint $(n,m)$ is
$\min\{n(r-1),|mA|-1\}=d_A(n,m)$, proving the slope assertion.
\end{proof}

\subsection{Collision orders and a finite tree formula}
The determinant form has a consequence not requiring interpolation
coordinates: it classifies arbitrary orders of tangency to a collision.
It also gives a finite optimization rule for all resolution phases.

\begin{theorem}[Collision-tree determination of memory phases]\label{thm:collision-tree}
Let $a(\theta)\in K$, $0\le\theta\le\theta_0$, be a known path.
For a fixed $m$, identify formal sums that are identical along the path.
Suppose that for every remaining pair of labels
\[
 |x_\alpha(\theta)-x_\beta(\theta)|
       \asymp\theta^{w_{\alpha\beta}},\qquad
 w_{\alpha\beta}\in[0,\infty),\quad 0<\theta\le\theta_0.
\]
After decreasing $\theta_0$ this hypothesis holds for every real analytic
path whose nonidentical sums are analytic at zero. Define
\begin{equation}\label{eq:tree-energy}
 \beta_{m,\ell}=
   \min_{\substack{J\text{ a set of remaining labels}\\|J|=\ell}}
         \sum_{\{\alpha,\beta\}\subset J}w_{\alpha\beta},
 \qquad \beta_{m,1}=0.
\end{equation}
Then $\mathcal V_{m,\ell}(a(\theta))\asymp
\theta^{\beta_{m,\ell}}$. For unavailable cardinalities the volume is
zero. Consequently the checkpoint law is
\begin{equation}\label{eq:tree-law}
 \overline R_{M;n,m}^{a(\theta)}\asymp
 \max_{\ell\le p_{n,m}}
       \theta^{2\beta_{m,\ell}/\ell}M^{-2/\ell},
\end{equation}
with zero-volume terms omitted; the streaming law is its maximum over
checkpoints. The exponents $\beta_{m,\ell}$ depend only on the pairwise
collision orders, and can be computed by a finite tree allocation rule.
\end{theorem}
\begin{proof}
Each fixed subset product has order equal to the sum of its pair orders.
There are finitely many subsets, so their maximum has the smallest such
order. This proves the determinant assertion and then
\eqref{eq:tree-law} by Theorem~\ref{thm:resolution-main}.
Constants may depend on the path's
nonzero leading coefficients; the intrinsic determinant theorem, in
contrast, is already uniform on $K$.

For completeness we give the allocation rule and its justification. The
triangle inequality for node differences implies
\[
 w_{\alpha\gamma}\ge
       \min\{w_{\alpha\beta},w_{\beta\gamma}\}.
\]
Indeed a smaller order on the left would contradict the upper bound by
the sum of the two differences on the right as $\theta\downarrow0$.
Thus thresholds on $w$ define nested equivalence classes. Make a rooted
tree from these classes, with root height zero; an internal cluster $C$
has height $h_C=\min_{\alpha\ne\beta\in C}w_{\alpha\beta}$,
and its children are the equivalence classes for orders strictly greater
than $h_C$. Split the root at orders greater than zero. Singleton
children are leaves. Redundant equal-height vertices can be omitted.
For an internal nonroot vertex put
$\Delta_C=h_C-h_{\operatorname{parent}(C)}$. The sum of these increments
on the common ancestor chain of two leaves is exactly their pair order.
Hence, for a selected leaf set $J$,
\begin{equation}\label{eq:tree-decomposition}
 \sum_{\{\alpha,\beta\}\subset J}w_{\alpha\beta}
   =\sum_{C\text{ internal},\ C\ne\mathrm{root}}
               \Delta_C\binom{|J\cap C|}{2}.
\end{equation}
Set $F_C(0)=F_C(1)=0$ for a leaf and $F_C(k)=+\infty$ for $k>1$.
For an internal vertex, including the root with $\Delta_{\rm root}=0$,
compute
\begin{equation}\label{eq:tree-dp}
 F_C(k)=\Delta_C\binom{k}{2}+
   \min_{\substack{\sum_{D\text{ child of }C}k_D=k}}
                 \sum_{D\text{ child of }C}F_D(k_D).
\end{equation}
Equation~\eqref{eq:tree-decomposition} proves this recurrence by induction.
It gives $F_{\rm root}(\ell)=\beta_{m,\ell}$. Successive finite
min-convolutions compute the table; no optimization over measures,
posterior encoders or codebooks is involved. For analytic paths the pair
orders are the finite orders of their first nonzero Taylor coefficients,
after identical sums have been identified. This proves the assertion.
\end{proof}

For the affine family of Section~\ref{sec:confluence}, different limiting
clusters have pair order zero, and distinct members of a cluster have pair
order one. Thus a choice of $k$ nodes from that cluster costs
$\binom{k}{2}$; its successive costs are $0,1,\ldots,s-1$.
Minimizing their sum picks the smallest costs from all clusters. Therefore
$\beta_{m,\ell}=\nu_{m,1}+\cdots+\nu_{m,\ell}$. This recovers
Theorem~\ref{thm:attainable-checkpoint}, while its separate proof remains
available. Higher orders of contact between distinct branches require the
pair orders above and are not encoded by the affine multiplicity list.

\subsection{Intersecting collision loci and arbitrary tangency order}
Here is a two-parameter consequence with different ranks on intersecting
loci and arbitrarily deep resolution phases along tangent approaches.
It illustrates the geometry of the determinant theorem rather than
requiring a separate asymptotic analysis for each path.

\begin{corollary}[A two-parameter collision arrangement]\label{cor:two-parameter}
Let $A_{u,v}=\{0,1,2+u,3+v\}$, with
$|u|,|v|\le1/16$, and take the four cells
\[
 k_0(t)=\tfrac14+\tfrac1{16}(t+t^{2+u}+t^{3+v}),\qquad
 k_i(t)=\tfrac14-\tfrac1{16}t^{a_i},\quad 1\le i\le3.
\]
Use any fixed admitted command cube and a fixed full-support prior.
For total horizon five set
$\rho=\max\{|u|,|v|\}$ and
$\tau=\min\{|u|,|v-u|,|v-2u|\}$. Then, uniformly on the square,
\begin{equation}\label{eq:two-parameter-law}
 \mathfrak R_{M,5}^{u,v}\asymp
 \max\left\{M^{-1/3},\ \rho^{1/2}M^{-1/4},\
           (\rho^2\tau)^{2/9}M^{-2/9}\right\}
\end{equation}
for either streaming criterion. The exact peak dimension is nine off the
three collision lines, eight on any such line away from their intersection,
and six at $(0,0)$.

Along $u=\theta$, $v=\theta+\theta^k$ with an integer $k\ge2$,
the three regimes have respective orders
\[
 M^{-1/3},\qquad \theta^{1/2}M^{-1/4},\qquad
       \theta^{(4+2k)/9}M^{-2/9}.
\]
Their crossover budgets have orders $\theta^{-6}$ and
$\theta^{-(8k-2)}$, with regrets of orders $\theta^2$ and
$\theta^{2k}$, respectively. These are comparison orders, not exact
integer thresholds.
\end{corollary}
\begin{proof}
The cells sum to one, lie between $3/16$ and $7/16$, and their coefficient
matrix has rank four: its last three rows separate the three nonconstant
coefficients and the sum of its rows gives the constant. They therefore
satisfy the compact-chamber hypotheses.

At checkpoint $n=3,m=2$ the nine positive formal exponents are
\[
 1,\quad 2,\ 2+u,\quad 3+u,\ 3+v,\quad
 4+2u,\ 4+v,\quad 5+u+v,\quad 6+2v.
\]
They form six uniformly separated groups, with three within-group gaps
$|u|,|v-u|,|v-2u|$. Denote these gaps in decreasing order by
$g_1\ge g_2\ge g_3$. For $\ell\le6$ the determinant volumes
are comparable to one. For $\ell=6+j$, $1\le j\le3$, they are
comparable to $g_1\cdots g_j$: choose one exponent from every group
and a second from the $j$ groups with largest gaps for the lower bound.
Any selection of $6+j$ exponents must double at least $j$ groups;
all other factors are bounded, proving the upper bound. Fixed powers of
$H$ only change constants. Since $p_{3,2}=9$, its profile is
\[
 \max\{M^{-1/3},\ g_1^{2/7}M^{-2/7},\
              (g_1g_2)^{1/4}M^{-1/4},\
              (g_1g_2g_3)^{2/9}M^{-2/9}\}.
\]
The signed differences satisfy $(v-u)-(v-2u)=u$, so the largest
absolute gap is at most twice the second largest. Moreover
$g_1\asymp\rho$, hence $g_2\asymp\rho$ and $g_3=\tau$.
The seven-dimensional term is the geometric interpolation, with weights
$3/7$ and $4/7$, of the six- and eight-dimensional terms, up to uniform
constants. This proves the expression in \eqref{eq:two-parameter-law}
for the checkpoint, also on the collision lines.

At the other checkpoints, the attainable dimensions are bounded by three,
six and three, respectively. Since every normalized gap is at most one,
their squared covering errors are at most $CM^{-1/3}$. The causal
classification therefore gives the same entire-horizon profile and the
checkpoint-three lower bound supplies its converse. Counting distinct
two-fold exponents gives the three exact peak dimensions asserted.

Finally, on the stated path $\rho\asymp\theta$ and
$\tau\asymp\theta^k$. Substitution gives the three displayed terms.
Equality in order of the first two gives $M=\theta^{-6}$, at regret
$\theta^2$; equality in order of the last two gives
$M=\theta^{-(8k-2)}$, at regret $\theta^{2k}$. As $k>1$, these
budgets are strictly separated as $\theta\downarrow0$ and all three
regimes occur. The exponent $k$ records contact with a collision locus,
not an additional assumed observable jet.
\end{proof}
